\textbf{Draft model (SmolLM2-135M).} Table~\ref{tab:epochsBd} shows the draft's validation behaviour over three epochs (\DTrainHours{}\,h on the CPU). Validation loss keeps falling (0.220, 0.194, 0.191), but the selection rule (highest validation macro-F1) picks epoch~\DBestEpoch{} (0.267) because macro-F1 is higher there than at epochs 2 and 3 (both 0.241); the differences are a few points on 206 rows and are not meaningful beyond showing that the 135M model learns the answer format and the label prior and little else. This is acceptable for a draft, whose purpose is to propose the target's likely tokens, not to be accurate by itself, but it means that the acceptance rates of the CPU pair were not benchmarked in a stored table (see below) and should not be inferred from the GPU pair.

\textbf{What the laptop is used for.} Configuration~B provides the model of the locally hosted application, whose \emph{Compare decoding} tab measures all four decoding modes live on the user's CPU (Section~\ref{sec:appres}); we did not run the full stored benchmark grid of Section~\ref{sec:kvres}--\ref{sec:specres} on the CPU, because the GPU grid already covers prompt lengths, output lengths, attention shapes and draft settings with a larger, more realistic model and any CPU table would add timings but no new correctness evidence. CPU-specific numbers should therefore be taken from the live tab rather than from this paper.

\begin{table}[t]\centering\caption{Config.~B draft (SmolLM2-135M + LoRA, CPU): validation after each epoch ($^*$ = selected).}\label{tab:epochsBd}
\small\begin{tabular}{rrrrr}\toprule
Epoch & Val. loss & Val. acc. & Val. macro-F1 & Val. evid. F1 \\\midrule
1$^*$ & 0.220 & 0.636 & 0.267 & 0.175 \\
2 & 0.194 & 0.631 & 0.241 & 0.048 \\
3 & 0.191 & 0.631 & 0.241 & 0.254 \\
\bottomrule\end{tabular}
\end{table}
